\section{Software architecture}\label{sec:architecture}

\statspai{} is implemented in \proglang{Python} on top of
\pkg{NumPy} \citep{harris2020array}, \pkg{SciPy}
\citep{virtanen2020scipy}, \pkg{pandas} \citep{mckinney2010data},
\pkg{statsmodels} \citep{seabold2010statsmodels}, \pkg{scikit-learn}
\citep{pedregosa2011scikit}, and \pkg{linearmodels}, with optional
backends kept behind extras and lazy imports. The architecture has four layers:
registry, dispatcher, result object, and output pipeline.

\subsection{Design principles}\label{sec:design-principles}

Seven principles govern the public API. A single flat namespace exposes
routine calls through \code{sp.<function>}. Mature estimator paths return
structured result objects rather than raw tuples, while auxiliary
helpers -- visualization and data-inspection utilities such as
\code{sp.panel\_view}, which returns Matplotlib handles plus an
\code{info} dictionary rather than an estimator result -- remain
registered with explicit stability and validation
metadata. Method families use dispatchers
(\code{sp.synth(method=...)}, \code{sp.dml(model=...)},
\code{sp.did(...)}). Reference parity is
attempted before broadening a validated claim. Failures are loud:
assumption violations raise typed exceptions or warnings and record
diagnostics. Output-changing changes go through the changelog and
\code{MIGRATION.md}. Finally, every public function must be
discoverable through the same registry used by human help functions and
agent schemas.

\subsection{Why an integration package?}\label{sec:why-integration}

A reasonable alternative is to compose \pkg{statsmodels},
\pkg{linearmodels}, \pkg{DoubleML}, \pkg{EconML}, \pkg{DoWhy}, and
method-specific \proglang{R} or \proglang{Stata} bridges directly,
which is often right for a single specialized model. The friction
appears when an applied project moves across identification designs:
diagnostics, citations, validation status, error handling, export
formats, and agent schemas no longer share one contract. \statspai{}
therefore treats wrapping as insufficient unless a method enters the
common registry, result-object surface, typed-error envelope, and
evidence ledger. Specialized packages remain the reference
implementations; the contribution here is to make cross-family handoffs
auditable and reproducible. The package is organized by identification
layer rather than by the historical boundaries of \proglang{R} and
\proglang{Stata} packages -- each layer reports its own design
diagnostics (balance and overlap, first-stage strength, cohort support
and event-study paths, bandwidth and density checks, donor weights and
placebos, cross-fit diagnostics) -- and the common point is not shared
identification assumptions but a shared post-estimation contract.

\subsection{Registry and discovery layer}\label{sec:registry}

The central data structure is a registry of \code{FunctionSpec}
records. Each entry carries a canonical function path, category, tags,
stability and validation status, parameter metadata, result class,
reference key, example call, and optional agent-card fields such as
assumptions, preconditions, failure modes, alternatives, limitations,
and typical sample-size guidance. The registry powers
\code{sp.help()}, \code{sp.list\_functions()},
\code{sp.search\_functions()}, \code{sp.describe\_function()},
\code{sp.function\_schema()}, \code{sp.agent\_card()}, the import-free
\code{schemas/} bundle, and the MCP tool catalogue. Tests assert that
exported public callables are registered and that validated/certified
entries expose their evidence status.

The validation tier is therefore not a paper-only artifact: it is
queryable from the running software for \emph{every} registered
function. \code{sp.help()} prints a per-tier census and points the user
to \code{sp.list\_functions(validation\_status="certified")} for the
validated core, while \code{sp.describe\_function(name)} prints a
\code{Stability}/\code{Validation} header with the attached evidence
notes before the description, so a user or agent landing on an
\code{api\_stable} or \code{experimental} helper sees its trust tier
before relying on its output.

\begin{lstlisting}[caption={The validation tier is visible through the
same discovery surface used by users and agents; this excerpt is
generated from \code{sp.describe\_function("callaway\_santanna")}.},%
label={lst:validation-tier-output},captionpos=b]
>>> spec = sp.describe_function("callaway_santanna")
>>> spec["stability"], spec["validation_status"]
('stable', 'certified')
>>> for note in spec["validation_notes"][:3]:
...     print(note)
Track A parity seed
R parity module 04_csdid
Stata parity module 04_csdid
>>> import textwrap
>>> print(textwrap.fill(spec["limitations"][0], width=64))
clustervars is not yet supported with bstrap=False; the analytic
standard errors do not account for within-cluster dependence, so
the multiplier bootstrap is required
\end{lstlisting}

\subsection{Family dispatchers}\label{sec:dispatcher}

Family dispatchers reduce method sprawl without hiding assumptions.
\code{sp.synth} dispatches classical, SDID, augmented, generalized,
matrix-completion, penalized, and related synthetic-control variants.
\code{sp.decompose} covers Blinder--Oaxaca, Gelbach, Fairlie, RIF,
DFL, Machado--Mata, Melly, counterfactual-distribution, and related
decompositions. \code{sp.dml} dispatches partially-linear,
partially-linear IV, interactive, and LATE-style DML variants. The
dispatcher contract is explicit: allowed \code{method=} values are
enumerated in the registry schema, and validated or mature variants
return compatible result surfaces so that summary, plotting, audit,
citation, and export calls do not change downstream.

The same uniformity extends to inference. Across the regression family
(\code{regress}, \code{feols}, \code{ivreg}, \code{panel},
\code{hdfe\_ols}), one canonical keyword, \code{vce=}, spells the full
standard-error menu: classical, heteroskedasticity-robust, one- and
two-way clustering, bias-reduced CR2/CR3
\citep{pustejovsky2018small}, the wild cluster bootstrap
\citep{cameron2008bootstrap, roodman2019fast}, and Conley spatial HAC
\citep{conley1999estimation}; fixed-effects Poisson and GLM
additionally expose CR2/CR3 and the restricted score wild bootstrap.
Each option is pinned to an external reference implementation ---
\proglang{Stata}'s \code{reghdfe}, \code{boottest}, \code{acreg}, and
\code{ppmlhdfe} \citep{correia2020fast}, and \proglang{R}'s
\pkg{clubSandwich} and
\pkg{sandwich} \citep{zeileis2020sandwich} --- and a coverage matrix in the package's continuous
integration ratchets the count of natively wired, externally validated
cells monotonically upward.

\subsection{Result-object hierarchy}\label{sec:results}

Most estimators return either \code{EconometricResults} or
\code{CausalResult}; specialized classes must implement the same public
surface. The standard contract covers summaries, plotting, dictionary
and JSON payloads, LaTeX/DOCX export, citations, audits, diagnostics,
and degradation records. This is why a mixed OLS/IV/DML/forest
table can be passed to \code{sp.regtable} (or its \proglang{R}-style
alias \code{sp.modelsummary}), and
why an MCP agent can keep only a result handle rather than copying
coefficient and covariance arrays between tool calls.

\subsection{Performance backbone}\label{sec:perfbackbone}

The performance-critical path is high-dimensional fixed-effect
absorption for fixed-effect and panel estimators. The submitted artifact
uses the reproducible \proglang{Python} path built on \pkg{NumPy} and
\pkg{Numba} \citep{lam2015numba}. An experimental Rust crate and selected \pkg{JAX} and
\pkg{PyTorch} paths exist behind extras, but the numerical and timing
claims in this paper do not depend on a compiled native extension.
Optional acceleration is treated as a backend choice, not as a separate
estimator or a hidden requirement.

\subsection{Output pipeline}\label{sec:output}

Publication output is unified through \code{sp.regtable},
\code{sp.modelsummary}, and result-level exporters. Mixed
\code{EconometricResults} and \code{CausalResult} objects can be
rendered as LaTeX, Markdown, HTML, Word, Excel, or JSON-safe agent
payloads. \code{sp.outreg2} gives a Stata-like migration surface, while
\code{sp.paper()} uses the same rendering layer to assemble tables,
figures, citations, and a report.

\subsection{The four-stage workflow}\label{sec:workflow}

The layers combine into a four-stage workflow -- detect
(\code{sp.detect\_design}, \code{sp.preflight}, \code{sp.recommend}),
fit (the estimator dispatchers), inspect (\code{summary()},
\code{plot()}, \code{sp.audit()}, sensitivity wrappers), and report
(\code{sp.regtable}, export methods, \code{cite()},
\code{sp.paper()}) -- with the same
boundaries available to human users and to MCP agents.
